\section{Fazit}
\label{sec:Fazit}

Die Pilotphase von GRACE befindet ist mit einer großen Herausforderung verbunden: Konfigurationswissen ist über mehrere Systeme fragmentiert verteilt. Die Ist-Analyse zeigt erhebliche Defizite. Die Vollständigkeit der Initialkonfigurationen beträgt lediglich 60-70\%. Durchschnittlich treten 10-15 Fehler pro Konfiguration auf. Die Durchlaufzeit beträgt ca. 8 Wochen. Es existiert keine externalisierte Wissensquelle. Wissensträgerabhängigkeit besteht bei wenigen Experten. Validierungsmechanismen besitzen keine strukturierte Form. Unter diesen erschwerten Bedingungen entwickelt die Arbeit eine Drei-Artefakte-Architektur. Diese externalisiert implizites Expertenwissen. Die Architektur umfasst drei Komponenten: GRACE Onboarding Fragebogen (96 → 15 Kernfragen), Entity Mapping Workbook (JSON-Schemas für 6 Entitätstypen durch Reverse Engineering) und Standard Operating Procedures.

Trotz der schwierigen Ausgangslage erreicht der aus den Artefakten entwickelte Master Creator im kontrollierten Experiment (n=18) beeindruckende Verbesserungen. Die Zeitreduktion beträgt 47\% (8:03 Min → 4:18 Min). Die Fehlerreduktion beträgt 81\% (5,00 → 0,94 Fehler). Die Vollständigkeit steigt von 86,1\% auf 98,1\%. Die zweistufige Verbesserung zeigt die Wirkungszusammenhänge deutlich. Strukturierung alleine führt zu einem Zuwachs von 16-26 Prozentpunkten (von 60-70\% auf 86,1\%). LLM-Automatisierung verstärkt die Wirkung um weitere 12 Prozentpunkte auf 98,1\%. Das zur Entwicklung genutzte 7B-Modell erfüllt Datensouveränitätsanforderungen (DSGVO, NDA). Es erreicht trotz Hardware-Limitationen referenzielle Integrität über 6 Entitätstypen. Die strategische Flexibilität (Self-Service, iterative Nutzung, Importer-Integration) ermöglicht Anpassung an heterogene Kundenanforderungen.

Die Arbeit unterliegt methodischen und technischen Einschränkungen. Der Fragebogen definiert gleichzeitig das Maß für Vollständigkeit. Dies kann zu zirkulärer Metrik führen. Die Qualität der Informationen ist jedoch durch systematische Externalisierung objektiv gestiegen. Vormals implizites Wissen liegt nun strukturiert dokumentiert vor. Technische Limitationen bestehen in der ID-Format-Diskrepanz zwischen Tool (kebab-case) und GRACE-Standard (UUID). Der Master Creator fokussiert zudem auf 6 Kern-Entitäten. Komplexere Strukturen wie itemAssociationRules oder variable Prozessformeln sind durch smartes Prompting generierbar. Diese sind jedoch noch nicht vollständig integriert. Die genannten Einschränkungen schmälern nicht den Kernbeitrag. Die Arbeit validiert erfolgreich LLM-gestützte Konfigurationsansätze für APS-Systeme unter realen Produktionsbedingungen. Sie demonstriert die Notwendigkeit systematischer Wissensexternalisierung als Enabler für LLM-Automatisierung.
